\documentclass[11pt]{article}
\usepackage[margin=1in]{geometry}
\usepackage{xcolor}
\usepackage{tcolorbox}
\usepackage{hyperref}
\usepackage{enumitem}
\setlist{noitemsep,topsep=0pt}
\usepackage{booktabs}
\usepackage{array}
\usepackage{amsmath}

% Define OSU orange and AI blue colors
\definecolor{osaorange}{RGB}{220,115,38}
\definecolor{aiblue}{RGB}{30,144,255}
\definecolor{correctgreen}{RGB}{34,139,34}

% Configure hyperref colors
\hypersetup{
    colorlinks=true,
    linkcolor=blue,
    urlcolor=blue,
    citecolor=blue
}

\title{}
\author{}
\date{}

\begin{document}

% Orange header box with black outline
\begin{tcolorbox}[colback=osaorange,colframe=black,boxrule=2pt,arc=0pt,left=6pt,right=6pt,top=10pt,bottom=10pt]
\centering
\textcolor{white}{\textbf{\Large Week 3 In-Class Worksheet - ANSWER KEY}}\\
\textcolor{white}{\textbf{\large Probability Distributions and Sampling Distributions}}
\end{tcolorbox}

\vspace{8pt}

\noindent\textbf{Course:} H524 Introduction to Biostatistics\\
\textbf{Topic:} Binomial, Normal, Sampling Distributions, and Confidence Intervals

\vspace{12pt}

\section{Problem 1: Binomial Distribution}

A new antibiotic has an 80\% success rate in treating a particular bacterial infection. A clinic treats 15 patients with this antibiotic.

\textbf{Part A:} What is the probability that exactly 12 patients are successfully treated?

\textbf{Given:} $n = 15$, $p = 0.80$, find $\Pr(X = 12)$

\textbf{Formula:} $\Pr(X = k) = \binom{n}{k} p^k (1-p)^{n-k}$

\textcolor{correctgreen}{\textbf{Solution:}}
\begin{align*}
\Pr(X = 12) &= \binom{15}{12} (0.80)^{12} (0.20)^{3}\\
&= \frac{15!}{12!3!} (0.80)^{12} (0.20)^{3}\\
&= 455 \times 0.06872 \times 0.008\\
&= \mathbf{0.2501}
\end{align*}

\vspace{0.5cm}

\textbf{Part B:} What is the probability that at least 13 patients are successfully treated?

Hint: $\Pr(X \geq 13) = \Pr(X=13) + \Pr(X=14) + \Pr(X=15)$

\textcolor{correctgreen}{\textbf{Solution:}}
\begin{align*}
\Pr(X=13) &= \binom{15}{13}(0.80)^{13}(0.20)^{2} = 0.2309\\
\Pr(X=14) &= \binom{15}{14}(0.80)^{14}(0.20)^{1} = 0.1319\\
\Pr(X=15) &= \binom{15}{15}(0.80)^{15}(0.20)^{0} = 0.0352\\
\Pr(X \geq 13) &= 0.2309 + 0.1319 + 0.0352 = \mathbf{0.3980}
\end{align*}

\vspace{0.5cm}

\textbf{Part C:} Calculate the expected number of successful treatments and the standard deviation.

Expected value: $\mu = np$

\textcolor{correctgreen}{\textbf{Solution:}} $\mu = 15 \times 0.80 = \mathbf{12}$

Standard deviation: $\sigma = \sqrt{np(1-p)}$

\textcolor{correctgreen}{\textbf{Solution:}} $\sigma = \sqrt{15 \times 0.80 \times 0.20} = \sqrt{2.4} = \mathbf{1.549}$

\vspace{0.5cm}

\textbf{Part D:} Use R to verify your answers. Write the R code you would use:

\textcolor{correctgreen}{\textbf{Solution:}}
\begin{verbatim}
# Part A: P(X = 12)
dbinom(12, size=15, prob=0.80)

# Part B: P(X >= 13)
sum(dbinom(13:15, 15, 0.80))
# OR: pbinom(12, 15, 0.80, lower.tail=FALSE)

# Part C: Expected value and SD
n <- 15
p <- 0.80
expected <- n * p
std_dev <- sqrt(n * p * (1-p))
\end{verbatim}

\newpage

\section{Problem 2: Normal Distribution}

Adult systolic blood pressure is normally distributed with mean 125 mmHg and standard deviation 18 mmHg.

\textbf{Part A:} What proportion of adults have blood pressure above 140~mmHg (hypertension)?

\textbf{Step 1:} Standardize using $Z = \frac{X - \mu}{\sigma}$

\textcolor{correctgreen}{\textbf{Solution:}}
$$Z = \frac{140 - 125}{18} = \frac{15}{18} = 0.833$$

\textbf{Step 2:} Find $\Pr(Z > z)$ using the standard normal table or R

\textcolor{correctgreen}{\textbf{Solution:}}
$$\Pr(Z > 0.833) = 1 - \Pr(Z < 0.833) = 1 - 0.7976 = \mathbf{0.2024}$$

About \textbf{20.2\%} of adults have blood pressure above 140 mmHg.

\vspace{0.5cm}

\textbf{Part B:} What proportion have blood pressure between 110 and 140 mmHg?

Hint: $\Pr(110 < X < 140) = \Pr(Z < z_2) - \Pr(Z < z_1)$

\textcolor{correctgreen}{\textbf{Solution:}}
\begin{align*}
Z_1 &= \frac{110 - 125}{18} = -0.833\\
Z_2 &= \frac{140 - 125}{18} = 0.833\\
\Pr(110 < X < 140) &= \Pr(Z < 0.833) - \Pr(Z < -0.833)\\
&= 0.7976 - 0.2024 = \mathbf{0.5952}
\end{align*}

About \textbf{59.5\%} have blood pressure in this range.

\vspace{0.5cm}

\textbf{Part C:} Find the blood pressure value that represents the 90th percentile.

\textbf{Step 1:} Find $z$ such that $\Pr(Z < z) = 0.90$

\textcolor{correctgreen}{\textbf{Solution:}} From Z-table or R: $z_{0.90} = 1.282$

\textbf{Step 2:} Convert back to original scale: $X = \mu + z\sigma$

\textcolor{correctgreen}{\textbf{Solution:}}
$$X = 125 + 1.282(18) = 125 + 23.08 = \mathbf{148.08 \text{ mmHg}}$$

\vspace{0.5cm}

\textbf{Part D:} Write the R code to solve parts A, B, and C:

\textcolor{correctgreen}{\textbf{Solution:}}
\begin{verbatim}
mu <- 125
sigma <- 18

# Part A: P(X > 140)
pnorm(140, mu, sigma, lower.tail=FALSE)

# Part B: P(110 < X < 140)
pnorm(140, mu, sigma) - pnorm(110, mu, sigma)

# Part C: 90th percentile
qnorm(0.90, mu, sigma)
\end{verbatim}

\newpage

\section{Problem 3: Sampling Distribution and Central Limit Theorem}

A population of cholesterol levels has mean $\mu = 200$ mg/dL and standard deviation $\sigma = 40$ mg/dL. A sample of $n = 64$ individuals is randomly selected.

\textbf{Part A:} What is the distribution of the sample mean $\bar{X}$?

By the Central Limit Theorem: $\bar{X} \sim N(\mu_{\bar{X}}, \sigma_{\bar{X}}^2)$

Mean of $\bar{X}$: $\mu_{\bar{X}} = $

\textcolor{correctgreen}{\textbf{Solution:}} $\mu_{\bar{X}} = \mu = \mathbf{200 \text{ mg/dL}}$

Standard error of $\bar{X}$: $\sigma_{\bar{X}} = \frac{\sigma}{\sqrt{n}} = $

\textcolor{correctgreen}{\textbf{Solution:}} $\sigma_{\bar{X}} = \frac{40}{\sqrt{64}} = \frac{40}{8} = \mathbf{5 \text{ mg/dL}}$

\textbf{Therefore:} $\bar{X} \sim N(200, 5^2)$ or $\bar{X} \sim N(200, 25)$

\vspace{0.5cm}

\textbf{Part B:} What is the probability that the sample mean is greater than 210 mg/dL?

\textbf{Step 1:} Standardize: $Z = \frac{\bar{X} - \mu_{\bar{X}}}{\sigma_{\bar{X}}}$

\textcolor{correctgreen}{\textbf{Solution:}}
$$Z = \frac{210 - 200}{5} = 2.0$$

\textbf{Step 2:} Find probability

\textcolor{correctgreen}{\textbf{Solution:}}
$$\Pr(\bar{X} > 210) = \Pr(Z > 2.0) = 1 - 0.9772 = \mathbf{0.0228}$$

R code: \texttt{pnorm(210, 200, 5, lower.tail=FALSE)}

\vspace{0.5cm}

\textbf{Part C:} What is the probability that the sample mean is between 195 and 205 mg/dL?

\textcolor{correctgreen}{\textbf{Solution:}}
\begin{align*}
Z_1 &= \frac{195 - 200}{5} = -1.0\\
Z_2 &= \frac{205 - 200}{5} = 1.0\\
\Pr(195 < \bar{X} < 205) &= \Pr(-1.0 < Z < 1.0)\\
&= \Pr(Z < 1.0) - \Pr(Z < -1.0)\\
&= 0.8413 - 0.1587 = \mathbf{0.6826}
\end{align*}

R code: \texttt{pnorm(205, 200, 5) - pnorm(195, 200, 5)}

\vspace{0.5cm}

\textbf{Part D:} How would your answers change if $n = 16$ instead of $n = 64$? Would that probability increase or decrease? Explain why.

\textcolor{correctgreen}{\textbf{Solution:}}

With $n = 16$: $\sigma_{\bar{X}} = \frac{40}{\sqrt{16}} = 10$ (larger SE)

That probability would \textbf{increase}. With a smaller sample size:
\begin{itemize}
\item Standard error is larger (10 vs. 5)
\item The sampling distribution is more spread out
\item Extreme values (like 210) are more likely
\item Z-score would be smaller: $Z = \frac{210-200}{10} = 1.0$ instead of 2.0
\item $\Pr(Z > 1.0) = 0.1587$ vs. $\Pr(Z > 2.0) = 0.0228$
\end{itemize}

\newpage

\section{Problem 4: Confidence Intervals}

A clinical trial measures the reduction in pain scores for 20 patients receiving a new medication. The sample mean reduction is $\bar{x} = 4.5$ points, with sample standard deviation $s = 1.8$ points.

\textbf{Part A:} Calculate a 95\% confidence interval for the true mean pain reduction.

\textbf{Step 1:} Determine what critical value to use. Should you use $z$ or $t$? Why?

\textcolor{correctgreen}{\textbf{Solution:}}

Use \textbf{t-distribution} because:
\begin{itemize}
\item Population SD ($\sigma$) is unknown
\item Small sample size ($n = 20 < 30$)
\end{itemize}

\textbf{Step 2:} Find the critical value.

Degrees of freedom = $n - 1 = $

\textcolor{correctgreen}{\textbf{Solution:}} $df = 20 - 1 = 19$

Critical value: $t_{0.975, df} = $ (use R: \texttt{qt(0.975, df)})

\textcolor{correctgreen}{\textbf{Solution:}} $t_{0.975, 19} = \mathbf{2.093}$

\textbf{Step 3:} Calculate the standard error: $\text{SE} = \frac{s}{\sqrt{n}}$

\textcolor{correctgreen}{\textbf{Solution:}}
$$SE = \frac{1.8}{\sqrt{20}} = \frac{1.8}{4.472} = \mathbf{0.403}$$

\textbf{Step 4:} Calculate the margin of error: $\text{ME} = t_{critical} \times \text{SE}$

\textcolor{correctgreen}{\textbf{Solution:}}
$$ME = 2.093 \times 0.403 = \mathbf{0.843}$$

\textbf{Step 5:} Construct the confidence interval: $\bar{x} \pm \text{ME}$

\textcolor{correctgreen}{\textbf{Solution:}}
$$CI = 4.5 \pm 0.843 = \mathbf{(3.66, 5.34) \text{ points}}$$

\vspace{0.5cm}

\textbf{Part B:} Interpret the confidence interval in context. Write a complete sentence explaining what this interval means.

\textcolor{correctgreen}{\textbf{Solution:}}

We are 95\% confident that the true mean pain reduction for all patients receiving this medication is between 3.66 and 5.34 points.

\vspace{0.5cm}

\textbf{Part C:} If you wanted a 99\% confidence interval instead, would it be wider or narrower than the 95\% interval? Explain.

\textcolor{correctgreen}{\textbf{Solution:}}

The 99\% CI would be \textbf{wider}. To be more confident that we've captured the true mean, we need a wider interval. The critical value increases ($t_{0.995, 19} = 2.861$ vs. $t_{0.975, 19} = 2.093$), resulting in a larger margin of error.

\newpage

\section{Problem 5: Integration Problem}

A pharmaceutical company is testing a new cholesterol-lowering drug. In the clinical trial:
\begin{itemize}
\item 100 patients receive the drug
\item Cholesterol reduction follows a normal distribution
\item Sample mean reduction: $\bar{x} = 28$ mg/dL
\item Sample standard deviation: $s = 12$ mg/dL
\end{itemize}

\textbf{Part A:} Calculate a 95\% confidence interval for the true mean cholesterol reduction.

\textcolor{correctgreen}{\textbf{Solution:}}

With large sample ($n=100$), use z-distribution: $z_{0.975} = 1.96$ (or t-distribution: $t_{0.975,99} \approx 1.984$)

\begin{align*}
SE &= \frac{s}{\sqrt{n}} = \frac{12}{\sqrt{100}} = 1.2\\
ME &= 1.96 \times 1.2 = 2.352\\
CI &= 28 \pm 2.352 = \mathbf{(25.65, 30.35) \text{ mg/dL}}
\end{align*}

\vspace{0.5cm}

\textbf{Part B:} Based on your confidence interval, is there strong evidence that the drug reduces cholesterol by more than 25 mg/dL on average? Explain your reasoning.

\textcolor{correctgreen}{\textbf{Solution:}}

\textbf{Yes}, there is strong evidence. The entire 95\% confidence interval (25.65, 30.35) is above 25 mg/dL. This means we can be 95\% confident that the true mean reduction exceeds 25 mg/dL.

\vspace{0.5cm}

\textbf{Part C:} If individual cholesterol reductions are normally distributed with the sample statistics above, approximately what proportion of patients will experience a reduction of at least 40 mg/dL?

Hint: Use the sample statistics as estimates of population parameters.

\textcolor{correctgreen}{\textbf{Solution:}}

Using $\mu = 28$ and $\sigma = 12$ as population estimates:

\begin{align*}
Z &= \frac{40 - 28}{12} = 1.0\\
\Pr(X \geq 40) &= \Pr(Z \geq 1.0) = 1 - 0.8413 = \mathbf{0.1587}
\end{align*}

About \textbf{15.9\%} of patients will experience a reduction of at least 40 mg/dL.

R code: \texttt{pnorm(40, 28, 12, lower.tail=FALSE)}

\vspace{0.5cm}

\textbf{Part D:} Explain the difference between the questions in Parts B and C. Why do we use different approaches (confidence interval vs. normal probability)?

\textcolor{correctgreen}{\textbf{Solution:}}

\textbf{Part B (Confidence Interval):} Asks about the \emph{population mean} - we use the sampling distribution of the sample mean, which has SE = $\sigma/\sqrt{n}$ = 1.2.

\textbf{Part C (Normal Probability):} Asks about \emph{individual patients} - we use the population distribution with SD = $\sigma$ = 12.

The key difference:
\begin{itemize}
\item Part B: Inference about a parameter (mean) $\rightarrow$ use SE
\item Part C: Prediction about individuals $\rightarrow$ use SD
\end{itemize}

Sample means are much less variable than individual observations, so the CI is narrower than the range containing 95\% of individuals.

\newpage

\section*{Reflection Questions}

\textbf{1.} What is the key difference between a standard deviation and a standard error?

\textcolor{correctgreen}{\textbf{Answer:}}

\textbf{Standard Deviation (SD):} Measures variability of individual observations around the mean. Describes the spread of the data.

\textbf{Standard Error (SE):} Measures variability of sample means around the population mean. Describes uncertainty in our estimate of the mean. $SE = \sigma/\sqrt{n}$, so it decreases with larger sample sizes.

\vspace{0.5cm}

\textbf{2.} When should you use a $t$-distribution instead of a normal distribution for confidence intervals?

\textcolor{correctgreen}{\textbf{Answer:}}

Use t-distribution when:
\begin{itemize}
\item Population standard deviation ($\sigma$) is unknown (must use sample SD $s$)
\item Sample size is small (typically $n < 30$)
\end{itemize}

With large samples ($n \geq 30$), the t-distribution closely approximates the normal distribution.

\vspace{0.5cm}

\textbf{3.} How does sample size affect the width of a confidence interval?

\textcolor{correctgreen}{\textbf{Answer:}}

As sample size increases, confidence intervals become \textbf{narrower} (more precise):
\begin{itemize}
\item Larger $n$ $\rightarrow$ smaller SE ($SE = s/\sqrt{n}$)
\item Smaller SE $\rightarrow$ smaller margin of error
\item Smaller ME $\rightarrow$ narrower CI
\end{itemize}

To cut the CI width in half, you need to quadruple the sample size.

\end{document}
